\subsection{A functional invariance principle for alternating BGW trees}
\label{ss:inv}

It is now time to consider specific offspring distributions.
These distributions will be chosen of the form
\begin{equation}
  \mubn(i)= \abn \cdot  (\bbn)^{i} \cdot \frac{(i+1)^{i-1}}{i!}  \quad (i \geq 0), \qquad \mucn(i)= \acn \cdot  (\bcn)^{i} \cdot \frac{(i+1)^{i-1}}{i!}  \quad (i \geq 0).
  \label{eq:mubn_mucn}
\end{equation}
Indeed, according to Proposition~\ref{prop:bitype},
alternating BGW trees with these distributions are connected to minimal factorizations
(for the moment, let us forget that, in Proposition~\ref{prop:bitype}, the root has a different offspring distribution).
We start by specifying the choices of the parameters
$\abn$, $\bbn$, $\acn$ and $\bcn$.
(Note that the chosen parameters and hence the offspring distributions depend on $n$.)

\begin{lemma}\label{lem:exist} For every $n \ge 1$,  fix $K_{n} \in \{1,\dots,n-1\}$.
\begin{enumerate}
  \item[(i)] We may choose positive parameters  $\abn,\bbn,\acn,\bcn$ in a unique way  such that \eqref{eq:mubn_mucn}
defines probability distributions with respective means
\[m^{n}_{\bullet}  \, \coloneqq  \, \sum_{i=0}^{\infty} i \cdot \mubn(i)= \frac{K_{n}+1}{n-K_{n}}, \qquad m^{n}_{\circ} \, \coloneqq \, \sum_{i=0}^{\infty} i \cdot \mucn(i)= \frac{n-K_{n}}{K_{n}+1}.\]
\item[(ii)] Assume that $K_{n} \rightarrow \infty$. If $\tfrac{K_{n}}{n} \rightarrow 0$ as $n \rightarrow \infty$, then
\[(\sbn)^{2} \sim \frac{K_{n}}{n}, \quad  \bbn \sim \frac{K_{n}}{n}  \qquad \text{ and } \qquad  (\scn)^{2} \sim   \big( \frac{n}{K_{n}} \big)^{3},\]
where $(\sbn)^{2}$ and $(\scn)^{2}$ denote respectively the variance of $\mubn$ and $\mucn$.
In particular, if $c>0$ is fixed and $\tfrac{K_{n}}{\sqrt{n}} \rightarrow c$ as $n \rightarrow \infty$, then   $ (\sbn)^{2} \sim \tfrac{c}{\sqrt{n}}$ and $(\scn)^{2} \sim  \frac{n^{3/2}}{c^{3}}$.
\end{enumerate}
\end{lemma}

\begin{proof}Let $F$ be the power series defined by
\begin{equation}
\label{eq:F}F(z)= \sum_{k \geq 0} \frac{(k+1)^{k-1}}{k!} z^{k}.
\end{equation}
It is elementary to check that $G(z)= \frac{z F'(z)}{F(z)}$ defines a continuous increasing function on $[0,1/e)$ and that $\lim_{z \rightarrow 0+} G(z)=0$, $\lim_{z \rightarrow 1/e-} G(z)=\infty$. One then chooses $\bbn, \bcn$ such that respectively $G(\bbn)=\frac{K_{n}+1}{n-K_{n}}$, $G(\bcn)= \frac{n-K_{n}}{K_{n}+1}$, and then $\abn=1/F(\bbn)$, $\acn=1/F(\bcn)$.
This proves (i).

For (ii), since $G(z)=z+o(z)$ as $z \rightarrow 0$ and since $G(\bbn)= \frac{K_{n}+1}{n-K_{n}}$, we get that $\bbn \sim \frac{K_{n}}{n}$ as $n \rightarrow \infty$.
The estimate for the variance $(\sbn)^{2}$ is then obtained by using the expression
\hbox{$(\sbn)^{2}= \frac{(\bbn)^{2} F''(\bbn)}{F(\bbn)}+\mbn-(\mbn)^{2}$:}
the dominant term when $n  \rightarrow \infty$ is $\mbn$
and we have \hbox{$(\sbn)^{2}\sim \frac{K_{n}}{n}$}.


The estimate concerning $(\scn)^{2}$ is more subtle, as it involves the behavior of $F$ near its radius of convergence $1/e$.
We can analytically extend $F$
on a complex split neighborhood on $1/e$ (i.e.\ on a complex neighborhood of $1/e$, without the real half-line $[1/e,+\infty)$).
Indeed we have $F(z)=-W(-z)/z$, where $W$ is the Lambert function (see \cite[Eq~3.1]{CGHJK96}).
Using \cite[Eq~4.22]{CGHJK96}, as $z \rightarrow 0$, we have
\begin{equation}
\label{eq:devF}F \left( \frac{1}{e}-z \right)=e- \sqrt{2} e^{3/2} \sqrt{z}+ o(\sqrt{z}),
\end{equation}
where $z$ is a complex number avoiding the negative real line
(throughout the paper, we use the principal determination of
$\sqrt{z}$ when $z$ is in $\CC \backslash \R_{<0}$).
By singular differentiation \cite[Theorem VI.8 p. 419]{FS09},
we get expansions for $F' \big( \tfrac{1}{e}-z \big)$ and $F'' \big( \tfrac{1}{e}-z \big)$
by differentiating the right-hand side of \eqref{eq:devF}.
Hence $G(1/e-z) \sim \frac{1}{\sqrt{2e}} \cdot \frac{1}{\sqrt{z}}$ as $z \rightarrow 0$.
By definition $\bcn$ is the solution of $G(\bcn)=\tfrac{n-K_n}{K_n+1}$, so that
\begin{equation}
  1/e-\bcn  \quad \mathop{\sim}_{n \rightarrow \infty} \quad  \frac{1}{2e} \left( \frac{K_n}{n-K_n}  \right)^{2}  \quad \mathop{\sim}_{n \rightarrow \infty} \quad  \frac{1}{2e} \left( \frac{K_n}{n}  \right)^{2}.
  \label{eq:est_bcirc}
\end{equation}
We again use the expression \hbox{$(\scn)^{2}= \frac{(\bcn)^{2} F''(\bcn)}{F(\bcn)}+m_{\circ}-m_{\circ}^{2}$}
to estimate the variance.
An easy computation gives $(\scn)^{2}\sim ( {n}/{K_{n}} )^{3}$.
(This time, the dominant term is $\left({(\bcn)^{2} F''(\bcn)}\right)\big/{F(\bcn)}$.)
\end{proof}
